\section{Introduction}
\label{sec:intro}

\subsection{Two Kinds of Set}

Call a set of natural numbers \emph{decidable} if some algorithm, given a
number, answers yes or no and always terminates.  Call it \emph{computably
enumerable}, or \ce., if some algorithm lists its members --- possibly
forever, in no particular order, with no promise about when any given
member appears.  Every decidable set is \ce.; the converse fails, and the
standard witness is the halting problem, the set of programs that
eventually stop.  One can list the halting programs by running all of them
a step at a time and reporting each as it stops.  One cannot decide
membership, because a program that has not stopped yet may stop later, and
nothing in the enumeration says which.

The asymmetry is the subject.  A \ce.\ set carries \emph{positive}
information that accumulates: a number may join the set at some stage of
the enumeration, and once it joins it never leaves.  Negative information
never becomes final.  Every construction in this paper is a way of
exploiting that asymmetry, and every difficulty in formalizing one comes
from the same place.

To compare sets we relativize.  Give an algorithm a free subroutine --- an
\emph{oracle} --- that answers membership questions about a set $B$.  If
some such algorithm decides $A$, write $A \leT B$: $A$ is \emph{Turing
reducible} to $B$, meaning $A$ is no harder than $B$.  A halting oracle
computation asks its oracle only finitely many questions, all below some
bound called the \emph{use} of that computation.  This finiteness is the
one lever the whole theory turns on, and we return to it below.

Reducibility sorts sets into layers of difficulty.  At the bottom sit the
decidable sets, reducible to everything.  Near the top sits the halting
problem, to which every \ce.\ set reduces.

\subsection{Post's Problem}

In 1944 Post asked what lies between \cite{post1944}.  Is every \ce.\ set
either decidable or as hard as the halting problem, or is there something
strictly in the middle?

The question resisted for twelve years.  It was settled in 1956--57, by
Friedberg and by Muchnik independently \cite{friedberg1957,muchnik1956},
and settled in a strong form: there are two \ce.\ sets $A$ and $B$ with
$A \not\leT B$ and $B \not\leT A$.  Neither can decide the other, so
neither is decidable and neither is as hard as the halting problem.  The
middle is not merely occupied; it is occupied by incomparable things.

\begin{figure}[h]
\centering
\begin{tikzpicture}[>=Stealth, every node/.style={font=\small}]
  \node (bot) at (0,0) {$\emptyset$ \ (the decidable sets)};
  \node (a) at (-2.2,1.7) {$A$};
  \node (b) at (2.2,1.7) {$B$};
  \node (top) at (0,3.4) {$K$ \ (the halting problem)};
  \draw (bot) -- (a);
  \draw (bot) -- (b);
  \draw (a) -- (top);
  \draw (b) -- (top);
  \node at (0,1.7) {\textit{no line}};
\end{tikzpicture}
\caption{The classical picture behind Post's problem: $A$ and $B$ sit
strictly between the decidable sets and the halting problem, and neither
reduces to the other.  The formalization reported here proves the two
missing edges --- $A\not\leT B$ and $B\not\leT A$ --- and proves that
neither set is decidable.  The halting problem itself does not appear in
the development, so the vertical placement in this diagram is classical
background rather than formalized content; see
Section~\ref{sec:scope}.}
\label{fig:degrees}
\end{figure}

A second classical result concerns how a \ce.\ set can be taken apart.
Sacks proved in 1963 that any \ce.\ set that is not decidable splits into
two disjoint \ce.\ pieces, neither of which can decide the whole
\cite{sacks1963}.  Nothing is thrown away --- the pieces reassemble to the
original --- and yet each piece has strictly less power than the set it
came from.  Both halves of the difficulty are genuinely present, and
neither half carries it alone.

\subsection{Why These Proofs Are Hard to Formalize}

Both results are proved by the same technique, introduced in the solution
of Post's problem and now called the \emph{priority method}.  It is worth
saying what the technique does before saying why it resists formalization.

To make $A \not\leT B$, one must defeat every candidate reduction.  Number
the oracle programs $\Phi_0, \Phi_1, \dots$ and impose one requirement per
program: $R_e$ says that $\Phi_e$ with oracle $B$ does not decide $A$.
Defeating a single $\Phi_e$ is easy in isolation.  Pick a number $x$ that
nothing else cares about --- a \emph{fresh witness} --- and keep it out of
$A$.  Watch $\Phi_e$ with oracle $B$ run on input $x$.  If it never
converges, it never decides $A$ and the requirement is met for free.  If it
converges and says ``$x \notin A$'', put $x$ into $A$.  Now the computation
is wrong, and the requirement is met by contradiction.

The difficulty is that the computation was run against $B$ \emph{as it
stood at that stage}, and $B$ is still being enumerated.  A later element
entering $B$ could change the computation and rescue it.  Here the finite
use is spent: the converging computation asked only finitely many questions
of $B$, so it suffices to \emph{restrain} $B$ below that use --- to promise
that no number smaller than the use will ever be enumerated into $B$.  The
computation is then frozen for good, and the diagonalization holds.

But restraints conflict.  A requirement that has promised to keep numbers
out of $B$ obstructs the requirements that want to put numbers in.  The
resolution is the priority ordering: requirements are ranked, and when a
higher-ranked one acts, every lower-ranked one is \emph{initialized} ---
its witness is discarded and its restraint cancelled, and it starts over
with a fresh witness chosen above the new restraint.  This is the
\emph{injury}.  The proof then argues that injury is finite: a requirement
is injured only by the finitely many stronger requirements above it, each
of which acts only finitely often, so eventually every requirement is left
alone forever and does its job.

This is a beautiful argument and an awkward one to make precise.  Its
content is carried by verbs in the future and past tense.  \emph{Choose a
witness nothing else cares about} --- nothing else meaning which
requirements, and cares in what sense, at which stage?  \emph{Wait until
the computation converges} --- waiting is not an operation; the
construction must at each stage make a decision from data it has, and
``converges'' is not such a datum, only ``converges within the resources
inspected so far''.  \emph{Restrain $B$ below the use} --- a promise about
all future stages, made at one stage, which some later stage must be shown
to keep.  \emph{Initialize everything weaker} --- an edit to unboundedly
many strategy records at once.  \emph{Eventually the stronger ones go
quiet} --- an induction along the priority order whose base case is that
the strongest requirement is never injured at all, and whose step needs a
measure of progress that no textbook states, because on paper it is
obvious.

In a proof assistant none of this can stay in the tense system.  Each
phrase has to become a piece of data or a theorem about it.  ``Fresh''
becomes a counter carried in the state and an invariant saying it dominates
everything mentioned so far.  ``Waiting'' becomes a decidable test on a
bounded execution, plus a separate theorem that a computation which really
does converge is eventually seen by that test.  ``Restraint'' becomes a
number in a record, plus a lemma that the enumeration respects it, plus the
preservation theorem that converts respect into a frozen computation.
``Initialize'' becomes an explicit rewriting of the strategy list.
``Eventually quiet'' becomes an induction along priority carrying an
explicit progress measure --- and finding that measure is real work, since
the informal proof supplies none.

That translation is the subject of this paper.  We report a Lean 4
formalization of both theorems --- Friedberg--Muchnik and Sacks splitting
in its splitting-with-non-reducibility form --- over one shared,
explicitly finitary model of oracle computation, and we report what the
translation cost and what it revealed.

\subsection{What This Paper Claims}
\label{sec:scope}

Our finding is not that the theorems are provable; they were proved
seventy years ago.  It is what happens when a second priority argument is
built on the first one's foundation.  The foundation transferred
completely, and one design decision made for the first theorem turned out
to be the load-bearing one for the second.  The priority layer transferred
not at all, and we argue that no generalization of it would have.  We
report that as a negative result about reuse: ``finite injury'' names a
family of arguments rather than a single reusable lemma.
Section~\ref{sec:formalization} is about both halves of that finding.

Two scope notes belong here rather than at the end.  First, what is
formalized is incomparability together with non-computability: for the two
constructed sets we prove $A \not\leT B$, $B \not\leT A$, and that neither
is decidable.  The classical step from there to ``strictly between
$\emptyset$ and the halting problem'' uses the fact that every \ce.\ set
reduces to the halting problem, and the halting problem, the jump, and
completeness do not appear in the development at all.  Figure~\ref{fig:degrees}
therefore shows more than we prove, and is labelled accordingly.  Second,
the full Sacks Splitting Theorem also makes the two halves \emph{low}; that
conclusion needs infinite-injury machinery and is deliberately out of
scope.  What we prove is the splitting-with-non-reducibility core.
